\begin{ouverture}[courbe]{8}{Fonctions : images et antécédents}{Partie 2 · Grandeurs et fonctions}
{Formules · Tableaux · Courbes · Programmes de calcul}
\objectif{Utiliser le vocabulaire de fonction.}
\objectif{Passer d’une formule à un tableau ou une courbe.}
\objectif{Modéliser une dépendance et en préciser les limites.}
\prerequis{Repérage et substitution dans une expression.}
\end{ouverture}
\pagestyle{otMain}\essentiel\fontsize{11}{15}\selectfont
\essSection{1}{Une entrée et une image}
\begin{definition}Une fonction associe à chaque nombre autorisé une seule image. On écrit $f(x)$ pour l’image de $x$ par $f$. Si $f(3)=8$, 8 est l’image de 3 et 3 est un antécédent de 8. Un nombre peut avoir plusieurs antécédents, ou aucun.\end{definition}
\begin{exemple}Pour $f(x)=x^2+1$, $f(-2)=5$. Les nombres $-2$ et 2 sont deux antécédents de 5. Aucun réel n’a pour image 0, car $x^2+1\ge1$.\end{exemple}
\begin{savoirfaire}[Lire un tableau]La première ligne contient les entrées, la seconde leurs images. Un tableau partiel ne donne pas toutes les valeurs d’une fonction : ne pas inventer l’image d’une entrée non indiquée.\end{savoirfaire}
\begin{definition}La courbe de $f$ rassemble les points de coordonnées $(x;f(x))$. L’abscisse est l’entrée ; l’ordonnée est son image. Vérifier l’appartenance d’un point exige de comparer son ordonnée à l’image de son abscisse.\end{definition}
\begin{attention}$f(2)$ ne signifie pas $f\times2$. Ne pas échanger l’image et l’antécédent dans une phrase.\end{attention}
\clearpage
\essSection{2}{Lire une courbe et construire un modèle}
\begin{savoirfaire}[Déterminer une image]Partir d’une abscisse, rejoindre la courbe verticalement, puis lire l’ordonnée. Pour chercher les antécédents d’un nombre, partir de cette ordonnée et repérer toutes les intersections horizontales avec la courbe.\end{savoirfaire}
\begin{attention}Une lecture graphique est souvent approchée. Une donnée exacte ou une formule peut permettre de justifier une valeur exacte ; la précision du dessin ne suffit pas toujours.\end{attention}
\begin{savoirfaire}[Modéliser]Nommer la variable et son unité, exprimer la quantité cherchée, puis préciser les entrées possibles. Un nombre de billets est entier ; une durée peut prendre des valeurs intermédiaires ; une longueur doit rester positive.\end{savoirfaire}
\begin{exemple}Pour une cuve de base d’aire 12 m$^2$, le volume en m$^3$ est $V(h)=12h$, tant que la hauteur $h$, en mètres, reste entre 0 et la hauteur maximale de la cuve.\end{exemple}
\begin{propriete}Comparer deux fonctions graphiquement revient à comparer leurs ordonnées pour une même abscisse. Les intersections correspondent aux entrées donnant la même image.\end{propriete}
\clearpage\rubrique{Découverte}
\begin{blocLivre}{ex-8-1}\begin{exercice}[1]{Calculer des images}
Pour $f(x)=3x-4$, calculer $f(0)$, $f(2)$ et $f(-3)$.
\end{exercice}\end{blocLivre}
\begin{blocLivre}{ex-8-2}\begin{exercice}[1]{Lire un tableau}
À l’aide du tableau, donner $f(-2)$ et les antécédents de 1 qui y figurent.\begin{center}\renewcommand{\arraystretch}{1.35}\begin{tabular}{c|cccc}\toprule
$x$ & $-2$ & $0$ & $1$ & $3$\\\midrule
$f(x)$ & $5$ & $1$ & $-1$ & $1$\\
\bottomrule\end{tabular}\end{center}
\end{exercice}\end{blocLivre}
\begin{blocLivre}{ex-8-3}\begin{exercice}[1]{Traduire une phrase}
Écrire avec la notation fonctionnelle : « 7 est l’image de $-1$ par $g$ ». Dire ensuite ce que signifie $g(4)=7$.
\end{exercice}\end{blocLivre}
\begin{blocLivre}{ex-8-4}\begin{exercice}[1]{Un point sur une courbe}
Pour $f(x)=x^2-2$, les points $A(3;7)$ et $B(-2;1)$ appartiennent-ils à la courbe ?
\end{exercice}\end{blocLivre}
\begin{blocLivre}{ex-8-5}\begin{exercice}[1]{Un antécédent algébrique}
Pour $h(x)=5x+2$, déterminer l’antécédent de 17.
\end{exercice}\end{blocLivre}

\clearpage\rubrique{Entraînement}
\begin{blocLivre}{ex-8-6}\begin{exercice}[2]{Une formule et un programme}
On choisit un nombre, on ajoute 3, puis on élève le résultat au carré. Donner la fonction obtenue et les images de 1 et $-3$.
\end{exercice}\end{blocLivre}
\begin{blocLivre}{ex-8-7}\begin{exercice}[2]{Lire une ligne brisée}
La courbe est formée des segments dessinés, avec les sommets sur les graduations indiquées. Lire $f(-1)$ et tous les antécédents de 1.
\label{fig-ex-8-7}\figureGeo{fonction-brisee}
\end{exercice}\end{blocLivre}
\begin{blocLivre}{ex-8-8}\begin{exercice}[2]{Plusieurs ou aucun antécédent}
Pour $f(x)=x^2$, déterminer les antécédents de 16, de 0 et de $-1$.
\label{fig-ex-8-8}\figureGeo{fonction-carre}
\end{exercice}\end{blocLivre}
\begin{blocLivre}{ex-8-9}\begin{exercice}[2]{Comparer deux représentations}
Les fonctions $f(x)=x+2$ et $g(x)=2x-1$ sont représentées. Pour quelle valeur de $x$ leurs images sont-elles égales ? Vérifier par un calcul.
\label{fig-ex-8-9}\figureGeo{fonction-intersection}
\end{exercice}\end{blocLivre}
\begin{blocLivre}{ex-8-10}\begin{exercice}[2]{Ce que ne dit pas un tableau}
Le tableau précédent permet-il de connaître exactement $f(2)$ ? Peut-on conclure que la fonction est une droite ?\begin{center}\renewcommand{\arraystretch}{1.35}\begin{tabular}{c|cccc}\toprule
$x$ & $-2$ & $0$ & $1$ & $3$\\\midrule
$f(x)$ & $5$ & $1$ & $-1$ & $1$\\
\bottomrule\end{tabular}\end{center}
\end{exercice}\end{blocLivre}

\clearpage\rubrique{Problèmes}
\begin{blocLivre}{ex-8-11}\begin{exercice}[3]{Un réservoir}
Un réservoir cylindrique de base d’aire 0,8 m$^2$ a une hauteur maximale de 1,5 m. Exprimer son volume en fonction de la hauteur $h$, puis trouver la hauteur correspondant à 600 L.
\label{fig-ex-8-11}\figureGeo{cuve-hauteur}
\end{exercice}\end{blocLivre}
\begin{blocLivre}{ex-8-12}\begin{exercice}[3]{Un tarif plafonné}
Un parking facture 2 € par heure avec un maximum de 10 € par journée. Les durées sont ici des nombres entiers d’heures entre 0 et 12. Donner les prix pour 3, 5 et 8 heures. La fonction est-elle linéaire sur toute la journée ?
\end{exercice}\end{blocLivre}
\begin{blocLivre}{ex-8-13}\begin{exercice}[3]{Des aires possibles}
Un rectangle de périmètre 40 cm a un côté de longueur $x$. Exprimer son aire. Quelles sont les entrées possibles ? Donner les images de 3 et de 17.
\label{fig-ex-8-13}\figureGeo{rectangle-fonction}
\end{exercice}\end{blocLivre}
\begin{blocLivre}{ex-8-14}\begin{exercice}[3]{Une lecture de déplacement}
La courbe donne la distance au point de départ pendant une promenade. Les segments ont pour sommets $(0;0)$, $(10;600)$, $(15;600)$ et $(25;0)$, en minutes et mètres. Décrire les phases et calculer la distance totale parcourue sur cet aller-retour rectiligne.
\label{fig-ex-8-14}\figureGeo{promenade}
\end{exercice}\end{blocLivre}
\begin{blocLivre}{ex-8-15}\begin{exercice}[3]{Résoudre par un produit nul}
Pour $f(x)=(x-2)(3x+6)$, déterminer les antécédents de 0. Donner ensuite une image strictement négative.
\end{exercice}\end{blocLivre}

\clearpage
